\subsection{VHDL} \label{sec:vhdl}
	VHDL (no angļu \termEn{VHSIC hardware description language}) ir aparatūras apraksta
	valoda, kura ieviesta 1980-ajos gados Amerikas Savienoto Valstu
	Aizsardzības ministrijas
	VHSIC (\termEn{very-high-speed integrated circuits})
	projekta ietvaros \cite[141.~lpp.]{VHSIC}\cite[1.~lpp.]{Perry-VHDL}.

	VHDL sintakse ir modelēta pēc Ada programmēšanas valodas sintakses, kur
	valodas konstrukcijas un to semantiskā nozīme pielāgota aparatūras aprakstam.
	VHDL sintakse tiek uzskatīta par visai izsmeļošu, t.i.,~valodas 
	konstrukcijas un izteiksmes ir pierakstāmas visai gari, kas seko no
	Ada sintakses, kā arī fakta, ka VHDL sākotnēji bija paredzēta kā 
	formāla valoda komponenšu dokumentācijai, kas viennozīmīgi aprakstītu
	komponenšu darbību.
	VHDL apraksti papildināja komponenšu dokumentāciju,
	aizstājot komplicēto un, bieži vien,
	nekonkrēto vai neviennozīmīgo rakstveidā dokumentēto darbības modeli.
	
	Tā kā VHDL apraksti seko noteiktai striktai, viennozīmīgai sintaksei,
	tie ir mašīnlasāmi un tika izstrādāta programmatūra, kas deva
	iespēju veikt aprakstītās shēmas simulāciju un vēlāk arī sintēzi.

	\begin{lstlisting}[language={[qucs]VHDL},label=kb:vhdl-example,gobble=4,%
			caption={VHDL apraksts loģiskajam \texttt{UN} elementam.}]
		-- This adds IEEE library and nine-valued logic
		library ieee;
		use ieee.std_logic_1164.all;
		
		entity andgate is
			port ( A : in std_logic;
			       B : in std_logic;
			       Q : out std_logic );
		end entity andgate;
		
		architecture rtl of andgate is
		begin
			Q <= A and B; -- Simply 'AND' the inputs and output result
		end architecture rtl;
	\end{lstlisting}

	Vienkāršs loģiskā elementa \texttt{UN} VHDL apraksta paraugs redzams
	\ref{kb:vhdl-example}~pirmkoda izdrukā. Piemērā redzami pamata sintaktiskās
	konstrukcijas komponentes aprakstam.
	Pirmkārt, entītija definē aprakstāmās komponentes ,,no ārpuses''
	redzamās daļas, no kurām svarīgākā (un piemērā vienīgā) ir pieslēgvietu
	saraksts; otrkārt, arhitektūra ir komponentes implementācijas definīcija,
	kurā tiek aprakstīta entītijā definētās komponentes darbība.

	Tā kā izmantojot aprakstīto komponenti hierarhiālā shēmas projektēšanā, 
	galvenokārt ir svarīga komponentes darbība, nevis tās konkrētā implementācija,
	VHDL apraksta arhitektūru var uzskatīt par ,,slēpto'' vai ,,privāto''
	komponentes daļu.

	VHDL iebūvētais \texttt{bit} tips (kas pieņem tikai |0| vai |1| vērtības)
	nevar parādīt tādus aparatūras loģikas stāvokļus, kā augstas impedances
	(|Z|) stāvoklis, atsaites rezistoru ,,vājos'' loģikas līmeņus un nezināmu
	(|X|) loģisko stāvokli. Šo problēmu atrisina IEEE 1164 standarts, kurš definē
	\texttt{std\_logic} tipu ar deviņiem dažādiem stāvokļiem:
	\begin{itemize}
		\item \verb|'U'| --- neinicializēta (nezināma) loģiskā vērtība;
		\item \verb|'X'| --- nezināma loģiskā vērtība no spēcīga devēja;
		\item \verb|'0'| --- loģiskā 0 no spēcīga devēja;
		\item \verb|'1'| --- loģiskais 1 no spēcīga devēja;
		\item \verb|'Z'| --- augstas impedances stāvoklis (nav devēja);
		\item \verb|'W'| --- nezināma loģiskā vērtība no vāja devēja;
		\item \verb|'L'| --- loģiskā 0 no vāja devēja (piem.,~atsaites rezistors pret zemi);
		\item \verb|'H'| --- loģiskais 1 no vāja devēja (piem.,~atsaites rezistors pret barošanas spriegumu);
		\item \verb|'-'| --- ignorējama vērtība (\termEn{don't care})%
			\footnote{Vērtību \verb|'-'| var izmantot salīdzināšanā,
				lai ignorētu konkrētus bitus no bitu vektora.},
	\end{itemize}
	kā arī definē operācijas ar \texttt{std\_logic} tipa signāliem vai mainīgajiem. 
	Praktiski visi VHDL simulācijas rīki šo bibliotēku piedāvā.
	
	Sintēzes rīki arī atbalsta \texttt{std\_logic} tipu, bet ne visas
	\texttt{std\_logic} tipa pieņemamās vērtības ir sintēzei saistošas.
	Reālu ciparu shēmu ieejas signāls vienmēr tiks interpretēts kā |0| (zems līmenis) 
	vai |1| (augsts līmenis), un izejas var pieņemt |0|, |1| vai augstas
	impedances (|Z|) stāvokli. Izvadīt nezināmu (|X|) loģisko vērtību
	fiziskai komponentei nav iespējams.

\subsection{Verilog} \label{sec:verilog}
\lipsum[5-6]

\clearpage
\subsection{Valodu salīdzinājums} \label{sec:hdl-comparison}
\lipsum[7-9]

	\begin{table}[thb]\small
		\centering
		\caption{VHDL un Verilog īpašību salīdzinājuma kopsavilkums.}
		\label{tbl:hdl-comparison}
		\begin{tabular}{p{0.2\textwidth}p{0.36\textwidth}p{0.36\textwidth}}
			\toprule
			~ & \textbf{VHDL} & \textbf{Verilog}\\ 
			\midrule
			\textbf{Sākotnējais\newline pielietojums} &
				Formāla komponenšu\newline dokumentēšana & Shēmu simulācija\\
			\midrule
			\textbf{Datu tipu\newline konversija} &
				Tikai ar speciālām konversijas funkcijām & 
				Vairumam tipu iespējama pēc\newline noklusējuma\\
			\midrule
			\textbf{Vairāku\newline stāvokļu loģika} &
				Ir, ar IEEE 1164 bibliotēku & Ir, iebūvēta\\
			\midrule
			\textbf{Pakotņu atbalsts} &
				Ir & Nav\\
			\midrule
			\textbf{Lietotāja definēti tipi} &
				Iespējami & Nav iespējami\\
			\midrule
			\textbf{Tranzistoru\newline modeļi} &
				Nav & Ir\\
			\midrule
			\textbf{Advancēti signālaizturu modeļi} &
				Ir (ar VITAL), bet nav domāti lietotāja izmantošanai &
				Ir, iebūvēti\\
			\bottomrule
		\end{tabular}
	\end{table}
	
	Kā redzams VHDL un Verilog valodu atšķirības ir daudz komplicētākas
	nekā tikai shēmas modeļu pieraksts, un, atkarībā no projekta specifikas,
	katrai valodai eksistē objektīvi iemesli tās izmantošanai. Īss pārskats
	valodu īpašību salīdzinājumam redzams \ref{tbl:hdl-comparison}~tabulā.
	
	Autora izvēlētā HDL mikrokontroliera kodola izstrādei ir VHDL. Kā viens
	no objektīviem iemesliem ir VHDL augsta līmeņa abstrakcijas, kas
	izstrādājot mikrokontroliera kodolu var ievērojami uzlabot
	produktivitāti. Darbā tiek izmantotas paša definētas
	pakotnes un datu tipi, un parametriski ģenerētās struktūras. Šādām
	konstrukcijām Verilog nav analogu.
	
	
